
\section{Preliminary Experiments}

The preliminary experiments were conducted in three stages: Initial sensor + window size experiment, Magnetometer feature experiment, and the final best performing model selection.

\subsection{Methodology}

In the preliminary stage, windows are grouped by recording session so that overlapping windows from a journey remain in the same partition. Different sessions from one participant may occur in both calibration (training) and holdout (testing) dataset. Models are trained on NOR-TMD or US-TMD, together with selected Sydney calibration sessions and evaluated on the remaining Sydney sessions.

The experiments utilise Optuna framework, an open-source Python framework used to automate model training and testing. For hyperparameter tuning, Tree-structured Parzen Estimator (TPE) algorithm is used, which is a type of Bayesian optimisation algorithm. Bayesian optimisation utilises previous evaluation scores to build a relatively inexpensive probability model to predict hyperparameter changes that would improve the model's performance. The classifier family candidates are a multilayer perceptron (MLP), Random Forest, and XGBoost. TPE evaluates and chooses family-specific hyperparameters, and a Hyperband pruner may stop a trial between grouped cross-validation folds. Selection uses macro F1 over five folds; only the overall winner of each search is retrained and evaluated on the session holdout. Completed-trial allocations are therefore unequal between families, and family-specific holdout superiority cannot be inferred from this design. For each run, 45 trials Bayesian optimisation budget is provided.

The three inertial sensors (accelerometer, gyroscope and magnetometer) records sensor values in 3-dimensional vector (x, y, z) form. To remove dependency on device orientation, each sensor value was reduced to a single scalar value using Euclidean magnitude $m=\sqrt{x^2+y^2+z^2}$ formula.

Figure~\ref{fig:preliminary-architecture} summarises the preliminary data preparation, feature extraction, and classifier evaluation procedure.

\begin{figure}[htbp]
\centering
\includegraphics[width=0.90\linewidth,height=0.65\textheight,keepaspectratio]{images/all-tmd-architecture.drawio.png}
\caption{Overview of the preliminary transport mode classification procedure. The testing-data connection denotes final evaluation of the selected classifier; hyperparameter selection uses training-fold data.}
\label{fig:preliminary-architecture}
\end{figure}

\subsection{Initial sensor + Window Size Experiment}

The initial sensor and window size experiment compared different sensor and window size configurations to explore how the changes influence the model performance and to establish a baseline configuration. This experiment used US-TMD and Sydney dataset together as training dataset.

The local experiment did not archive an immutable manifest of the complete collector dataset at execution time. Consequently, historical totals for collected sessions, participants, duration, and raw samples could not be reconstructed reliably. The eight non-pressure trials used 63 collector session groups, divided into 31 calibration and 32 holdout groups. The four pressure-enabled trials used 47 session groups, divided into 23 calibration and 24 holdout groups. Consequently, comparisons involving pressure are confounded by the change in the effective collector cohort and should not be interpreted as measuring the effect of pressure alone. The description of the dataset cohort is in Table~\ref{tab:transfer-v2-effective-cohorts}.

\begin{table}[htbp]
\centering
\caption{Effective collector cohorts}
\label{tab:transfer-v2-effective-cohorts}
\begin{threeparttable}
\begin{tabular}{lrrrr}
\hline
\textbf{Sensors}\tnote{1} & \textbf{Runs} & \textbf{Calibration} &
\textbf{Holdout} & \textbf{Total session groups} \\
\hline
A, G       & 4 & 31 & 32 & 63 \\
A, G, M    & 4 & 31 & 32 & 63 \\
A, G, M, B & 4 & 23 & 24 & 47 \\
\hline
\end{tabular}
\begin{tablenotes}
\item[1] A: Accelermoeter, G: Gyroscope, M: Magnetometer, B: Barometer
\end{tablenotes}
\end{threeparttable}
\end{table}

\subsubsection{Results}

% A = accelerometer, G = gyroscope, M = magnetometer, P = pressure.
\begin{table}[htbp]
\centering
\caption{Macro F1 for each window and sensor combination.}
\label{tab:transfer-v2-run-f1}
\begin{threeparttable}
\begin{tabular}{|c|p{2cm}p{2cm}p{2cm}|c|}
\hline
\textbf{Window (s)}\tnote{1} & A, G & A, G, M & A, G, M, B & Average \\
\hline
10 & 0.5766 & 0.6267 & 0.7118 & 0.6383 \\
20 & 0.6189 & 0.6789 & 0.7376 & 0.6785 \\
30 & 0.6409 & 0.6772 & 0.7837 & 0.7006 \\
60 & 0.6328 & 0.7255 & 0.7480 & 0.7021 \\
\hline
Average & 0.6173 & 0.6771 & 0.7453 & 0.6799 \\
\hline
\end{tabular}
\begin{tablenotes}
\item[1] A: Accelermoeter, G: Gyroscope, M: Magnetometer, B: Barometer
\end{tablenotes}
\end{threeparttable}
\end{table}

% Collector holdout metrics from the local ALL-TMD Transfer v2 MLflow experiment.
% Pressure-enabled runs use a smaller effective collector cohort.
\begin{table}[htbp]
\centering
\caption{Collector holdout accuracy and balanced accuracy for each initial sensor and window size run.}
\label{tab:transfer-v2-run-accuracy}
\begin{threeparttable}
\small
\begin{tabular}{llrr}
\hline
\textbf{Window (s)} & \textbf{Sensors}\tnote{1} & \textbf{Accuracy} & \textbf{Balanced accuracy} \\
\hline
10 & A, G & 0.7458 & 0.6408 \\
10 & A, G, M & 0.8419 & 0.6527 \\
10 & A, G, M, B & 0.8529 & 0.8213 \\
20 & A, G & 0.7821 & 0.7211 \\
20 & A, G, M & 0.8901 & 0.6961 \\
20 & A, G, M, B & 0.8949 & 0.7867 \\
30 & A, G & 0.8060 & 0.7421 \\
30 & A, G, M & 0.9042 & 0.6835 \\
30 & A, G, M, B & 0.9276 & 0.7932 \\
60 & A, G & 0.7786 & 0.7665 \\
60 & A, G, M & 0.8998 & 0.7861 \\
60 & A, G, M, B & 0.9063 & 0.7994 \\
\hline
\end{tabular}
\begin{tablenotes}
\item[1] A: Accelerometer, G: Gyroscope, M: Magnetometer, B: Barometer
\end{tablenotes}
\end{threeparttable}
\end{table}

Adding the magnetometer and barometer to the base sensor set (accerlerometer and gyroscope) displayed increase in classification performance.  All four trials that used pressure outperformed the eight trials without it, but the pressure-enabled trials used a smaller subset of session groups because built-in barometer was not available for all devices. The apparent gain therefore confounds the effect of pressure data with the inherent difference in training and testing dataset.

The increase in window length also monotonic increase in average macro F1 score. This was anticipated because longer windows can suppress transient noise and provide more observations for stable statistical features. However, 30-second pressure run had the highest score than that of 60-second pressure. This is likely due to small dataset size and data noise. The results are reported in Table~\ref{tab:transfer-v2-run-f1}.

\subsection{Pressure-feature and Four Transport Classes Experiment}

The pressure-feature and four transport classes experiment compared three pressure configurations: no pressure, pressure change from the window start, and pressure change from the session baseline. This experiment used NOR-TMD and Sydney dataset together as training dataset.

\subsubsection{Sydney Collector Data Snapshot}

Table~\ref{tab:aws-run-9c676a2c-effective-collector-snapshot-1} describes the effective collector session snapshot for four no-pressure trials. Table~\ref{tab:aws-run-9c676a2c-effective-collector-snapshot-2} describes the effective collector session snapshot for 8 pressure-enabled trials.

% Effective collector snapshot group: No-pressure trials.
% Effective collector snapshot source: unique raw-snapshot subset matching MLflow collector membership.
% Effective collector snapshot session ID digest: cc01e3571ae19ca8eb2061624b171af1e87bdf7d33502011096ed0656c04aacd.
\begin{table}[htbp]
\centering
\caption{Effective collector session snapshot for 4 no-pressure trials.}
\label{tab:aws-run-9c676a2c-effective-collector-snapshot-1}
\begin{tabular}{lrrrr}
\hline
\textbf{Mode} & \textbf{Sessions} & \textbf{Participants} & \textbf{Duration} & \textbf{Samples} \\
\hline
Bus & 15 & 4 & 05:47:17 & 1,119,558 \\
Car & 70 & 5 & 32:38:47 & 8,291,501 \\
Train & 46 & 4 & 20:46:34 & 4,119,463 \\
Tram & 9 & 1 & 02:04:44 & 371,846 \\
Total & 140 & 7 & 61:17:21 & 13,902,368 \\
\hline
\end{tabular}
\end{table}

% Effective collector snapshot group: Pressure-enabled trials.
% Effective collector snapshot source: unique raw-snapshot subset matching MLflow collector membership.
% Effective collector snapshot session ID digest: dc9e2299b57ebf6c1d859b5608bc738b55496eb4f822b6542f46510a3ebde38b.
\begin{table}[htbp]
\centering
\caption{Effective collector session snapshot for 8 pressure-enabled trials.}
\label{tab:aws-run-9c676a2c-effective-collector-snapshot-2}
\begin{tabular}{lrrrr}
\hline
\textbf{Mode} & \textbf{Sessions} & \textbf{Participants} & \textbf{Duration} & \textbf{Samples} \\
\hline
Bus & 12 & 2 & 05:28:14 & 1,062,233 \\
Car & 67 & 4 & 32:05:43 & 8,192,902 \\
Train & 25 & 2 & 16:28:17 & 3,349,854 \\
Total & 104 & 5 & 54:02:13 & 12,604,989 \\
\hline
\end{tabular}
\end{table}


\subsubsection{Results}

\begin{table}[htbp]
\centering
\caption{Macro F1 for each window and pressure configuration combination.}
\label{tab:aws-run-9c676a2c-trial-f1}
\begin{tabular}{lrrr}
\hline
\textbf{Window (s)} & \textbf{No pressure} & \textbf{window-start $\Delta$} & \textbf{session-baseline $\Delta$} \\
\hline
10 & 0.5603 & 0.5454 & 0.5367 \\
20 & 0.5870 & 0.5470 & 0.5316 \\
30 & 0.5996 & 0.5539 & 0.5412 \\
60 & \textbf{0.6053} & 0.5617 & 0.5543 \\
\hline
\end{tabular}
\end{table}

% Potential TODO add f1 score bar chart for each transport mode

% Collector holdout metrics from 20260809T055942Z-9c676a2c-full.
% Pressure-enabled trials have a different effective cohort and omit tram.
\begin{table}[htbp]
\centering
\caption{Collector holdout accuracy and balanced accuracy for each pressure-feature run.}
\label{tab:aws-run-9c676a2c-trial-accuracy}
\small
\begin{tabular}{llrr}
\hline
\textbf{Window (s)} & \textbf{Pressure configuration} & \textbf{Accuracy} & \textbf{Balanced accuracy} \\
\hline
10 & No pressure & 0.7789 & 0.6269 \\
20 & No pressure & 0.7783 & 0.6670 \\
30 & No pressure & 0.7899 & 0.6770 \\
60 & No pressure & 0.7972 & 0.6818 \\
10 & Window-start $\Delta$ & 0.7989 & 0.7291 \\
20 & Window-start $\Delta$ & 0.8000 & 0.7347 \\
30 & Window-start $\Delta$ & 0.8005 & 0.7576 \\
60 & Window-start $\Delta$ & 0.8193 & 0.7611 \\
10 & Session-baseline $\Delta$ & 0.7775 & 0.7279 \\
20 & Session-baseline $\Delta$ & 0.7723 & 0.7096 \\
30 & Session-baseline $\Delta$ & 0.7934 & 0.7185 \\
60 & Session-baseline $\Delta$ & 0.8019 & 0.7565 \\
\hline
\end{tabular}
\end{table}

Table~\ref{tab:aws-run-9c676a2c-trial-f1} shows the macro F1 scores of the pressure-feature and four transport classes experiment. The metrics demonstrate that all four trials that did not use pressure outperformed the eight trials which used pressure, but the pressure-enabled trials excluded tram because the only tram participant's device lacked a barometer. The apparent gain therefore confounds the effect of pressure with a change from four classes to three. Furthermore, a substantial amount of session data were collected with devices that did not include a built-in barometer were excluded from the eight pressure-enabled trials. Therefore, caution must be taken before attempting a head-to-head comparison.

\subsection{Magnetometer Feature Experiment}

The magnetometer feature experiment compared three magnetometer feature sets while holding the accelerometer and gyroscope features constant, without barometer features. This experiment used NOR-TMD and Sydney dataset together as training dataset.

\subsubsection{Sydney Collector Data Snapshot}

Table~\ref{tab:aws-run-a46cc6a0-effective-collector-snapshot} reports the effective Sydney collector data snapshot taken at the time of the experiment execution.

\begin{table}[htbp]
\centering
\caption{Effective collector session snapshot.}
\label{tab:aws-run-a46cc6a0-effective-collector-snapshot}
\begin{tabular}{lrrrr}
\hline
\textbf{Mode} & \textbf{Sessions} & \textbf{Participants} & \textbf{Duration} & \textbf{Samples} \\
\hline
Bus & 35 & 4 & 18:59:05 & 3,726,582 \\
Car & 71 & 5 & 33:00:03 & 8,356,739 \\
Train & 46 & 4 & 20:46:34 & 4,119,463 \\
\hline
Total & 152 & 7 & 72:45:42 & 16,202,784 \\
\hline
\end{tabular}
\end{table}

\subsubsection{Results}

The best feature combination used a 60-second window and magnetometer standard deviation, range, and minimum, achieving macro F1 score of 0.8804. Adding magnetometer minimum helped, while adding mean produced virtually no further gain. Longer windows consistently performed better. Table~\ref{tab:aws-run-a46cc6a0-trial-f1} shows the macro F1 scores of the magnetometer feature experiment.

\begin{table}[htbp]
\centering
\caption{Collector holdout macro F1 for each window and magnetometer features combination.}
\label{tab:aws-run-a46cc6a0-trial-f1}
\begin{threeparttable}
\begin{tabular}{|c|c|c|c|}
\hline
\textbf{Window (s)} & \textbf{SD\tnote{1} + range} & \textbf{+ minimum} & \textbf{+ mean} \\
\hline
10 & 0.8194 & 0.8160 & 0.8144 \\
20 & 0.8424 & 0.8435 & 0.8385 \\
30 & 0.8534 & 0.8543 & 0.8525 \\
60 & 0.8643 & \textbf{0.8804} & 0.8800 \\
\hline
\end{tabular}
\begin{tablenotes}
\item[1] SD: Standard deviation
\end{tablenotes}
\end{threeparttable}
\end{table}

% Collector holdout metrics from 20260812T125950Z-a46cc6a0-full.
\begin{table}[htbp]
\centering
\caption{Collector holdout accuracy and balanced accuracy for each magnetometer-feature run.}
\label{tab:aws-run-a46cc6a0-trial-accuracy}
\begin{threeparttable}
\small
\begin{tabular}{llrr}
\hline
\textbf{Window (s)} & \textbf{Magnetometer features} & \textbf{Accuracy} & \textbf{Balanced accuracy} \\
\hline
10 & SD\tnote{1} + range & 0.8193 & 0.8244 \\
20 & SD + range & 0.8413 & 0.8475 \\
30 & SD + range & 0.8516 & 0.8590 \\
60 & SD + range & 0.8630 & 0.8692 \\
10 & SD + range + minimum & 0.8177 & 0.8198 \\
20 & SD + range + minimum & 0.8438 & 0.8472 \\
30 & SD + range + minimum & 0.8533 & 0.8584 \\
60 & SD + range + minimum & 0.8787 & 0.8848 \\
10 & SD + range + minimum + mean & 0.8163 & 0.8181 \\
20 & SD + range + minimum + mean & 0.8393 & 0.8423 \\
30 & SD + range + minimum + mean & 0.8521 & 0.8560 \\
60 & SD + range + minimum + mean & 0.8792 & 0.8840 \\
\hline
\end{tabular}
\begin{tablenotes}
\item[1] SD: Standard deviation
\end{tablenotes}
\end{threeparttable}
\end{table}

\subsection{Calibration Fraction Experiment}

The calibration fraction experiment compared three different fractions of Sydney calibration dataset while holding the accelerometer, gyroscope, and
magnetometer features constant, without barometer features. This experiment used NOR-TMD and Sydney dataset together as training dataset.

\subsubsection{Sydney Collector Data Snapshot}

Table~\ref{tab:aws-run-c11f9a1a-effective-collector-snapshot} reports the effective Sydney collector data snapshot taken at the time of the experiment execution.

\begin{table}[htbp]
\centering
\caption{Effective collector session snapshot.}
\label{tab:aws-run-c11f9a1a-effective-collector-snapshot}
\begin{tabular}{lrrrr}
\hline
\textbf{Mode} & \textbf{Sessions} & \textbf{Participants} & \textbf{Duration} & \textbf{Samples} \\
\hline
Bus & 38 & 4 & 20:27:45 & 4,018,646 \\
Car & 71 & 5 & 33:00:03 & 8,356,739 \\
Train & 46 & 4 & 20:46:34 & 4,119,463 \\
Total & 155 & 7 & 74:14:21 & 16,494,848 \\
\hline
\end{tabular}
\end{table}

\subsubsection{Results}

\begin{table}[htbp]
\centering
\caption{Macro F1 with different calibration fraction}
\label{tab:transport-class-calibration-macro-f1}
\begin{threeparttable}
\small
\begin{tabular}{|l|c|ccc|c|}
\hline
\textbf{Condition} & \textbf{Calibration\tnote{1}} & \textbf{Bus} & \textbf{Car} & \textbf{Train} & \textbf{Average} \\
\hline
A & 0.8, 0.4, 0.4 & 0.6308 & 0.9378 & 0.9122 & 0.8269 \\
B & 0.4, 0.8, 0.4 & 0.8419 & 0.8259 & 0.9231 & 0.8637 \\
C & 0.4, 0.4, 0.8 & 0.8351 & 0.9251 & 0.8301 & 0.8634 \\
\hline
\end{tabular}
\begin{tablenotes}
\item[1] Proportions are ordered as bus, car, and train
\end{tablenotes}
\end{threeparttable}
\end{table}

% Collector holdout metrics from 20260814T114054Z-c11f9a1a-full.
\begin{table}[htbp]
\centering
\caption{Collector holdout accuracy and balanced accuracy for each calibration-fraction run.}
\label{tab:aws-run-c11f9a1a-trial-accuracy}
\small
\begin{tabular}{lrrr}
\hline
\textbf{Condition} & \textbf{Calibration (Bus, Car, Train)} & \textbf{Accuracy} & \textbf{Balanced accuracy} \\
\hline
A & 0.8, 0.4, 0.4 & 0.8899 & 0.8631 \\
B & 0.4, 0.8, 0.4 & 0.8724 & 0.8875 \\
C & 0.4, 0.4, 0.8 & 0.8846 & 0.8886 \\
\hline
\end{tabular}
\end{table}

\begin{figure}[htbp]
\centering
\includegraphics[width=0.8\linewidth,height=0.45\textheight,keepaspectratio]{report_segments/images/collector-holdout-class-f1-by-calibration-allocation.png}
\caption{Macro F1 with different calibration fraction}
\label{fig:f1-by-calibration-allocation}
\end{figure}

Table~\ref{tab:transport-class-calibration-macro-f1} and Figure~\ref{fig:f1-by-calibration-allocation} shows the individual and average macro F1 scores of the calibration fraction experiment. Changing the fractions also changes the composition and size of each testing set, therefore, these scores should be compared with some caution.


\subsection{Analysis}

For every preliminary experiments, the winning classifier of final search was XGBoost.

% The experimental pipeline for calibration fraction experiment already assigns inverse-frequency class weights during both cross-validation and final training, and it tunes using F1. This gives each class substantially equal influence despite different window counts. Therefore,


The four preliminary experiments should not be compared head-to-head since each experiment's dataset composition differs significantly, especially, the initial sensor and window size experiment used US-TMD as training dataset, whereas the other three experiments used NOR-TMD. Furthermore, the data composition of the three NOR-TMD experiments are not identical, since the experiments were conducted while the Sydney data collection was in progress. Moreover, pressure-feature experiment included tram in its classification classes.
